\begin{rmk}
注意由定义可知，\eqref{Sat isom} 中的两个环都有一组以 $\xcoch(T)^\sigma/W_0$ 为指标的自然基。即每个 $\mu\in \xcoch(T)^\sigma/W_0$ 给出特征函数 $1_{K\mu(\varpi)K}\in H_G$，以及函数 $\chi_{_{V_\mu}}\in \mathbf J$，其中 $V_\mu$ 是 ${^L}G$ 的「最高权 $\mu$」不可约表示。
然而，上例表明同构 $\Sat^{cl}$ 一般不把 $1_{K\mu(\varpi)K}$ 送到 $\chi_{_{V_\mu}}$。
\end{rmk}

\begin{rmk}
\label{R: norm Sat}
当然，可以定义同构 $\Sat'^{cl}: H_{\GL_2}\to R(\GL_2)$ 把 $T_p$ 送到 $\tr$（它在 $\bZ$ 上有定义，且与 $p^{1/2}$ 的选取无关）。这与 Satake 同构的归一化问题有关（见 \cite[\S 8]{Gr}）。更根本地，它与 Langlands 对偶群的\textbf{正确}定义有关。本文不讨论这一问题（注 \ref{R:Tate twist} 除外）。从几何 Satake 角度所作的详细讨论可参看 \cite[\S 5.5]{Zh2}。
\end{rmk}

\medskip

现在我们解释几何的做法：我们不再把 $G(F)$ 看作拓扑群、考虑其上 $K$-双不变的紧支函数空间，而是把 $G(F)$ 看作某个代数几何对象，研究其上的 $K$-双等变层范畴。等价而更方便的做法是，把 $G(F)/K$ 看作代数几何对象，研究其上的 $K$-等变层范畴。更确切地说，我们将构造 $k$ 上一个无穷维代数簇 $\Gr_G$，称为 $G$ 的仿射 Grassmann 簇，使 $G(F)/K$ 是它的 $k$-点集。
要看出这为何可行，我们先分析 $G(F)/K$ 的某些子集，它们可以自然地与某个代数簇的 $k$-点集等同。

\begin{ex}\label{minu Sch}
设 $G=\GL_n$ 定义在 $F$ 上。则商 $G(F)/K$ 可等同于 $F^n$ 中的格构成的集合。这里所谓格，指 $F^n$ 的有限自由 $\mO$-子模 $\Lambda$，且满足 $\Lambda\otimes_{\mO}F=F^n$。例如 $\La_0:=\mO^n\subset F^n$ 是格，通常称为标准格。双射由
\begin{equation}\label{coset vs lattice}
\GL_n(F)/K\simeq \big\{\mbox{格 } F^n\mbox{ 中的格}\big\},\quad gK\mapsto \Lambda:=g\La_0.
\end{equation}
给出。

考虑 $\omega_i=1^i0^{n-i}\in \bZ^{n}$，并如例 \ref{coweightGLn} 那样把它看作 $\GL_n$ 的余权。则容易看出，
在等同 \eqref{coset vs lattice} 下，集合 $K\omega_i(\varpi)K/K$ 可等同于满足
\[\big\{\varpi\mO^n\subset \Lambda \subset \mO^n,\quad  \dim_{k}\mO^n/\Lambda=i\big\}.\]
的格 $\Lambda\subset F^n$ 的集合。映射
$$\Lambda\mapsto (k^n=\mO^n/\varpi\mO^n\to \mO^n/\La)$$
把此集合等同于 $k^n$ 的 $i$ 维商空间（等价地，$k^n$ 的 $(n-i)$ 维子空间）的集合。总之，$K\omega_i(\varpi)K/K$ 与 Grassmann 簇 $\Gr(n-i,n)$（固定 $n$ 维空间中 $(n-i)$ 维平面的簇）的 $k$-点集之间有典范双射。自然期望 $\Gr(n-i,n)$ 是 $\Gr$ 的子簇。
\end{ex}

\begin{ex}\label{quasiminuSch}
设 $G=\GL_2$ 定义在 $F$ 上。依据 Satake 同构（特别见 \eqref{upper tr}），考虑集合
\begin{equation}\label{Gr2}
K\begin{pmatrix}\varpi^2 & \\ &1 \end{pmatrix}K/K\sqcup \begin{pmatrix}\varpi & \\ &\varpi \end{pmatrix} K/K.
\end{equation} 是合理的。在等同 \eqref{coset vs lattice} 下，它可等同于满足
\begin{equation}\label{Gr2latt}
\{\La\subset \mO^2\mid \on{length}\mO^2/\La=2\}.
\end{equation}
的格 $\La\subset F^2$ 的集合。
若 $F=k((\varpi))$，可把 $\mO^2/\La$ 看作 $k^4=\mO^2/\varpi^2\mO^2$ 的 $2$ 维商空间。此外，乘以 $\varpi$ 在 $\mO^2/\varpi^2\mO^2$ 上诱导幂零自同态，它保持 $\mO^2/\La$。容易看出，$\mO^2/\varpi^2\mO^2$ 的每个在此幂零自同态下稳定的 $2$ 维商空间都这样得到。因此 \eqref{Gr2} 可等同于 $\Gr(2,4)$ 的某个闭子簇的 $\bF_q$-点集。

若 $F=\bQ_p$，情形更复杂，因为我们不能把 $\mO^2/\La$ 看作 $k=\bF_p$ 上的向量空间。为克服这一困难，自然考虑另一个集合
\begin{equation}\label{wGr2latt}
\{\La\subset\La'\subset \mO^2\mid \dim\mO^2/\La'=\La/\La'=1\}.
\end{equation}
有一个自然满射 $m$ 把 \eqref{wGr2latt} 送到 \eqref{Gr2latt}，即忘掉 $\La'$。容易验证，除 $\La=p\La_0$ 之外，$m^{-1}(\La)$ 只含一个元素。
 
与前面类似的推理表明，\eqref{wGr2latt} 等同于 Hirzebruch 曲面 $\bP(\mO_{\bP^1}(-1)+\mO_{\bP^1}(1))$ 的 $\bF_p$-点集。此外，集合 $p^{-1}(\La_0)$ 构成该曲面上一条 $(-2)$-曲线的 $\bF_p$-点集。沿这条曲线收缩，所得射影簇的 $\bF_p$-点集自然等同于 \eqref{Gr2latt}。
\end{ex}

现在我们给出 $G$ 的仿射 Grassmann 簇的精确刻画。本节余下部分，我们允许 $F$ 稍更一般，即假定
$F$ 是关于离散赋值完备的局部域，整数环为 $\mO$，剩余域 $k$ 是\textbf{完全}的，特征 $p>0$\footnote{在等特征情形，即 $\cha F=\cha k$ 时，此假设并非必要。此处加上它是为了统一处理。}。设 $\varpi\in\mO$ 为一致化元。

回忆 $k$-代数 $R$ 称为完全的，指 Frobenius 自同态 $R\to R, \ r\mapsto r^p$ 是双射。
对完全 $k$-代数 $R$，令 $W(R)$ 为 $R$ 的（$p$-典型）Witt 向量环。它是元素形如序列 $(r_0,r_1,r_2,\ldots)\in R^{\bN}$ 的环，加法与乘法由某些（复杂的）多项式给出。到第一个分量的投影 $W(R)\to R, \ (r_0,r_1,\ldots)\mapsto r_0$ 是满环同态，带有乘性（但\textbf{非加性}）的截面 $R\to W(R), \ r\mapsto [r]=(r,0,0,\ldots)$，称为 $r$ 的 Teichm\"uller 提升。若 $R$ 是完全环，$W(R)$ 中每个元素可唯一地写成 $\sum_{i\geq 0} [a_i]p^i$，故 $W(R)$ 可视为「以 $p$ 为变量、系数在 $R$ 中的幂级数环」。
例如，$W(\bF_p)=\bZ_p$。我们定义 $R$ 上系数在 $\mO$ 中的 Witt 向量环为
\begin{equation}\label{ramified Witt vector}
W_\mO(R):= W(R)\hat\otimes_{W(k)}\mO:= \underleftarrow\lim_n W_{\mO,n}(R),\quad W_{\mO,n}(R):=W(R)\otimes_{W(k)} \mO/\varpi^n.
\end{equation}
特别地，若 $R=\bar k$，我们有时把 $W_\mO(\bar k)$ 记作 $\mO_L$，把 $W_{\mO}(\bar k)[1/\varpi]$ 记作 $L$。
注意若 $\cha F=\cha k$，则 $W_\mO(R)\simeq R[[\varpi]]$，且 $r\mapsto [r]$ 是环同态。我们有时写
\begin{equation}\label{discs}
D_{n,R}=\Spec W_{\mO,n}(R),\quad D_R=\Spec W_\mO(R),\quad D_R^*=\Spec W_\mO(R)[1/\varpi],
\end{equation}
并把它们看作由 $\Spec R$ 参数化的（去心）圆盘族。
给定 $D_R$ 上两个 $G$-torsor $\mE_1$ 与 $\mE_2$，从 $\mE_1$ 到 $\mE_2$ 的一个修正（modification）是同构 $\mE_1|_{D_R^*}\simeq \mE_2|_{D_R^*}$。我们通常记作 $\mE_1\dashrightarrow\mE_2$，以示该同构并非定义在整个 $D_R$ 上。

我们需要在完全代数几何中工作。一些基础可参看 \cite[附录 A]{Zh1}、\cite{BS} 与 \cite[\S A.1]{XZ}。
